\chapter*{Zakończenie}
\addcontentsline{toc}{chapter}{Zakończenie}

Celem pracy było opracowanie symulatora sieci LoRa Mesh opartych na projekcie Meshtastic oraz ocena wiarygodności jego wyników. Cel ten został zrealizowany. Opracowany symulator pozwala rozmieścić węzły na mapie, obliczyć łącza radiowe z uwzględnieniem ukształtowania terenu i zabudowy, uruchomić sieć złożoną z rzeczywistego oprogramowania węzłów oraz przeanalizować przebieg transmisji i retransmisji każdego pakietu.

Najważniejszą cechą symulatora jest rozdzielenie warstwy logicznej od warstwy radiowej. Routing, retransmisje i role węzłów realizowane są przez oprogramowanie \texttt{meshtasticd} uruchamiane w osobnym kontenerze dla każdego węzła, dzięki czemu zachowanie węzłów w symulacji odpowiada zachowaniu urządzeń pracujących w sieci, również po kolejnych zmianach protokołu. Warstwa radiowa łączy model Longley-Rice z danymi SRTM, stratę dyfrakcyjną od budynków z modeli LoD1, stałą stratę otoczenia, trwałe cieniowanie łączy oraz fading. Model współdzielenia kanału uwzględnia czas trwania transmisji, efekt przechwycenia, pracę w półdupleksie oraz emulację mechanizmu wykrywania zajętości kanału, którego moduł symulowanego radia w oprogramowaniu węzła nie wykonuje. Ślad symulacji zapisywany jest w bazie InfluxDB wraz z przyczyną każdej nieudanej próby odbioru, co pozwala oddzielić straty wynikające z zasięgu od strat wynikających ze współdzielenia kanału.

Weryfikacja przeprowadzona na danych z polskiej sieci Meshtastic wykazała, że:
\begin{itemize}
\item model łączy poprawnie porządkuje łącza od najlepszych do najgorszych (AUC równe 0,79 na danych niezależnych od kalibracji), a kalibracja straty otoczenia i cieniowania trzykrotnie zmniejsza błąd przewidywanych prawdopodobieństw (wynik Briera 0,17 wobec 0,51),
\item model skalibrowany odwzorowuje sposób przekazania pakietów, bezpośrednio lub przez pośrednika, w 75\% par węzłów, a w obszarze z danymi o zabudowie zgodność routingu z siecią jest wysoka,
\item trasy w symulacji są średnio o około 1,5 skoku krótsze niż w sieci rzeczywistej, co wynika z przeszacowania liczby łączy oraz z braku w symulacji węzłów, które nie podają dokładnej pozycji,
\item model współdzielenia kanału zachowuje się zgodnie z modelem teoretycznym protokołu ze sprawdzaniem kanału przy małym obciążeniu i daje wyniki znacznie lepsze niż protokół ALOHA,
\item symulator działa poprawnie dla sieci liczących do 80 węzłów, a czas uruchomienia scenariusza rośnie liniowo z liczbą węzłów.
\end{itemize}

Symulator nadaje się zatem przede wszystkim do porównywania wariantów sieci, takich jak rozmieszczenie węzłów, ich role, limity skoków czy presety modemu, oraz do badania zjawisk związanych ze współdzieleniem kanału i skalowalnością. Bezwzględne wartości uzyskane dla konkretnej istniejącej sieci należy interpretować ostrożnie i wyznaczać na podstawie serii przebiegów, ponieważ oprogramowanie węzłów wprowadza niedeterministyczność, której symulator nie kontroluje.

Symulator posiada również ograniczenia. Jedna, globalna wartość straty otoczenia nie odwzorowuje różnic pomiędzy instalacjami poszczególnych węzłów, takich jak umieszczenie anteny czy jakość przewodów. Model kolizji zakłada pracę wszystkich węzłów z jednym współczynnikiem rozpraszania i nie sumuje mocy wielu jednoczesnych zakłóceń. Nie jest modelowane ograniczenie cyklu pracy nadajnika (ang. \textit{duty cycle}) obowiązujące w paśmie 869,4--869,65 MHz, które w rzeczywistej sieci może ograniczać liczbę retransmisji przy dużym obciążeniu. Model zabudowy wymaga danych LoD1, które w ramach pracy zaimportowano wyłącznie dla okolic Rzeszowa. Skalowalność sprawdzono dla sieci liczących do 80 węzłów, a uruchamianie jednego kontenera na węzeł sprawia, że większe sieci wymagają odpowiednio większych zasobów komputera.

Kierunki dalszych prac wynikają bezpośrednio z tych ograniczeń. Najważniejszym z nich jest przeprowadzenie kontrolowanej kampanii pomiarowej z węzłami o znanej konfiguracji, która pozwoliłaby oddzielić błędy modelu propagacji od niepewności danych o istniejących węzłach. Model radiowy można rozszerzyć o ograniczenie cyklu pracy, sumowanie mocy zakłóceń oraz częściową ortogonalność różnych współczynników rozpraszania, co umożliwiłoby badanie sieci z węzłami pracującymi na różnych presetach. Rozproszenie kontenerów węzłów na wiele komputerów pozwoliłoby symulować sieci liczące kilkaset węzłów. Ponieważ warstwa radiowa symulatora jest niezależna od oprogramowania węzłów, możliwe jest również dołączenie innych rozwiązań, na przykład MeshCore, i porównanie ich działania w tych samych warunkach propagacyjnych.
